% 本文件由 LaTeX 排版 Agent 根据有效 translation.json 自动生成。
% author=Apocrypha; work_id=0806; title=救主童年阿拉伯语福音

\chapter{救主童年阿拉伯语福音}

% structure: p0001 source=h1 target=omitted reason=page-title-already-rendered-by-parent

在圣父、圣子及圣神之名，唯一天主。\par

赖至高者之助与恩宠，我们开始书写吾主、导师及救主耶稣基督的奇迹之书，即称为\WorkTitle{童年福音}者：在主之平安内。阿们。\par

1. 我们在那生活于基督时代的大司祭若瑟之书中发现如下记载。有人称该大司祭为盖法。他述说耶稣曾发言，且确实，当祂躺在摇篮中时，对其母玛利亚说：\Quote{我是耶稣，天主子，\OriginalTerm{圣言}{Logos}，即你所生育者，恰如加俾额尔天使向你预报的；我父派遣我来是为世界的救恩。}\par

2. 在亚历山大纪元第三百零九年，奥古斯都颁布谕令，命各人归于本籍登记。若瑟遂起身，携其配偶玛利亚往耶路撒冷，来到白冷，为与其家族在本籍登记。及至一山洞，玛利亚告知若瑟分娩之期已近，不能入城；却说：\Quote{我们且进此洞。}此事发生在日落时分。若瑟急忙出去，欲寻一妇人相助。当他在忙碌此事时，见一来自耶路撒冷的希伯来老妇，便说：\Quote{请来，好妇人，进此洞，洞中有一妇人即将分娩。}\par

3. 于是，日落之后，老妇与若瑟同至洞中，二人进入。看啊，洞中充满光明，比灯烛之光更美，比日光更灿烂。那婴孩裹着襁褓，正在哺乳圣母玛利亚之乳，置于马槽中。二人对此光惊叹，老妇问圣母玛利亚：\Quote{你是这孩子的母亲吗？}圣母玛利亚表示肯定，她便说：\Quote{你全然不像厄娃的女儿们。}圣母玛利亚说：\Quote{我的儿子在孩童中无与伦比，同样他的母亲在妇女中也无与伦比。}老妇回答：\Quote{我的主母，我来求报酬；我患瘫痪已久。}我们的主母圣母玛利亚对她说：\Quote{将你的手放在这孩子身上。}老妇照做，立即痊愈。于是她出去说：\Quote{从今以后，我愿终身做这孩子的侍从和仆人。}\par

4. 然后牧人们来到；他们点燃火堆，极其欢乐时，天上万军显现，赞美歌颂至高天主。牧人们也同样赞美时，那山洞当时竟成了如同上天的圣殿，因天上与地上的声音一起光荣赞美天主，由于主基督的诞生。那位希伯来老妇见到这些奇迹的彰显，感谢天主说：\Quote{我感谢你，天主，以色列的天主，因为我亲眼看见了救世主的诞生。}\par

5. 割损之期，即第八日，临近了，孩子当依法受割损。于是他们在洞中为祂行了割损。那希伯来老妇取了一块皮；但有人说是她取了脐带，置于一古纳尔多油瓶中。她有一子，是香料商，便将此瓶交给他，说：\Quote{你看，即使有人向你出价三百第纳尔，也切勿出售这瓶纳尔多油。}这正是罪妇玛利亚所购买，倾倒在吾主耶稣基督头上和脚上，随后用她的头发擦干的那瓶油。\BibleRef{路 7:37}十日后，他们将祂带到耶路撒冷；在祂出生后第四十日\BibleRef{肋 12:4}之后，他们将祂带进圣殿，置于主前，按梅瑟法律的诫命为祂献祭，即：\Quote{凡开胎首生的男子，当称为圣于主。}\par

6. 当时年迈的西默盎看见祂如光柱照耀，当时圣母玛利亚，祂的童贞母亲，怀抱祂，欢喜快乐。天使们赞美祂，如同侍卫环绕君王一般，立于祂周围。西默盎于是急忙上前，向圣母玛利亚伸出双手，对主基督说：\Quote{主啊，如今请照祢的话，让祢的仆人平安离去；因为我亲眼看见了祢的怜悯，即祢为万民的救恩所预备的，照亮万邦的光，祢子民以色列的荣耀。}女先知亚纳也在场，上前来感谢天主，称圣母玛利亚为有福的。\BibleRef{路 2:25}\par

7. 当主耶稣在犹太的白冷诞生，在黑落德王时代，看啊，有贤士从东方来到耶路撒冷，正如则拉杜士特所预言的；他们携带着礼物：黄金、乳香和没药。他们朝拜了祂，并向祂献上礼物。于是圣母玛利亚取了其中一条襁褓带，因手头拮据，便给了他们；他们以极大的尊敬接受。同时有天使化作先前引领他们旅途的那颗星显现；他们便追随星光前行，直到返回本国。\BibleRef{玛 2:1}\par

8. 他们的君王和首领聚集到他们那里，询问他们所看见或所行的事，以及他们如何往返、带来了什么。他们便将圣母玛利亚所给的那条襁褓布展示给他们看。于是他们举行庆节，按他们的习俗点燃火焰并崇拜之，并将那襁褓布投入火中；火焰抓住它并包裹了它。当火焰熄灭后，他们取出襁褓布，竟与先前一模一样，如同未被火烧过一般。因此他们开始亲吻它，将其置于头上和眼前，说：\Quote{这实在是确凿无疑的真理。确实，火焰不能烧毁它，这是大事。}于是他们怀着极大的尊敬将其收在宝库中。\par

9. 当黑落德见贤士离他而去，没有返回他那里，就召集司祭和智者，对他们说：\Quote{告诉我基督将在哪里诞生。}当他们回答说：\Quote{在犹太的白冷}，他就开始图谋杀害主耶稣基督。于是，主的一位天使在梦中显现给若瑟，说：\Quote{起来，带着孩子和祂的母亲，逃往埃及。}\BibleRef{玛 2:13} 他因此在鸡鸣时分起身，出发了。\par

10. 当他正思忖如何启程时，清晨降临在他身上，他只走了很短的一段路。此刻他正走近一座大城，城中有一尊偶像，埃及的其他偶像和神祇都向这尊偶像献供品和许愿。有一位司祭侍立在这尊偶像前，侍奉它；每当撒殚从这偶像中说话，他就向埃及及其领土的居民报告。这位司祭有一个三岁的儿子，被几个魔鬼缠身；他说许多话、发许多声音；当魔鬼抓住他时，他撕破衣服，赤身露体，向人们扔石头。那座城里有一所旅舍，是奉献给那偶像的。当若瑟和玛利亚来到那城，转入那旅舍时，市民非常害怕；所有首领和偶像的司祭都聚集到那偶像前，对它说：\Quote{我们土地上兴起的这骚动和混乱是什么？}偶像回答他们说：\Quote{有一位天主秘密地来到了这里，祂是真天主；除了祂，没有别的神堪受敬拜，因为祂确实是天主子。}当这土地意识到祂的临在时，因祂的到来而战栗、震动；我们因祂大能的伟大而极其恐惧。在同一时刻，那偶像倒塌了，所有埃及居民和其他人因它的倒塌而聚集起来。\par

11. 那司祭的儿子，旧病复发，进了旅舍，遇见了若瑟和玛利亚，其他人都已逃开。玛利亚洗了主基督的襁褓，晾在一些木头上。那附魔的男孩因此过来，拿起一块襁褓，放在自己头上。于是魔鬼以乌鸦和蛇的形状从他口中逃出。男孩因主基督的命令立即痊愈，开始赞美天主，然后感谢治好他的主。他父亲见他恢复健康，说：\Quote{我儿，你发生了什么事？你是怎样得医治的？}儿子回答：\Quote{当魔鬼把我摔倒在地时，我进了旅舍，在那里遇见一位庄严的妇人和一个男孩，她刚洗过的襁褓晾在一些木头上；我拿起其中一块，放在头上，魔鬼就离开我逃走了。}父亲因此大为欢喜，说：\Quote{我儿，这孩子可能就是永生天主子，祂创造了天地；因为祂来到我们这里时，偶像破碎了，所有神祇都倒塌了，因祂威严的大能而灭亡。}\par

12. 这里应验了那预言说：\BibleQuote{从埃及我召回了我的儿子。}的确，若瑟和玛利亚听说那偶像倒塌灭亡后，就战栗害怕。然后他们说：\Quote{我们在以色列地时，黑落德想要杀害耶稣，因此杀了白冷及其四境所有的婴孩；毫无疑问，埃及人一听说这偶像被毁，就会用火烧死我们。}\par

13. 他们从那里出发，来到一个地方，那里有强盗，他们抢劫了几个人的行李和衣服，并把他们捆绑起来。那时强盗听见巨大的响声，好像一位辉煌的君王率军、战车和鼓从城中出来；强盗因此惊恐，丢下所有掠夺物。被掳的人站起来，彼此解开捆绑，取回行李，走了。他们看见若瑟和玛利亚来到那地方，就对他们说：\Quote{那位君王在哪里？听到祂来临的辉煌响声，强盗就丢弃了我们，我们得以安全逃脱。}若瑟回答他们：\Quote{祂将跟在我们后面。}\par

14. 此后他们来到另一座城，那里有一个附魔的妇人，撒殚，这受诅咒的反叛者，在一次她夜间出去打水时缠上了她。她既不能穿衣服，也不能住在屋里；每当他们用锁链和皮条捆绑她，她就挣断，赤身逃到旷野；她站在十字路口和墓地，向人们扔石头，给她的亲友带来极重的灾祸。当玛利亚看见她，就怜悯她；撒殚立刻离开她，化作一个少年的形像逃走了，说：\Quote{玛利亚，你和你的儿子，使我遭殃！}那妇人就从折磨中被治愈，恢复理智，因赤身而害羞，避开人们的视线，回家到亲友那里。她穿上衣服后，将事情经过告诉了父亲和亲友；由于他们是城中的首领，就以极大的尊敬和款待接待了玛利亚和若瑟。\par

15. 次日，他们得到旅途所需之物后便离开了，当晚到达另一城镇，那里正在举行婚宴；但因受诅咒的撒殚的诡计和巫术的作为，新娘成了哑巴，不能说话。当玛利亚进入城镇，怀抱着她的儿子主基督时，那哑巴新娘看见她，就向主基督伸出双手，把祂抱过来，搂在怀里，亲吻祂，俯身摇动祂的身体。立刻她的舌结松开了，耳朵也开了；她感谢赞美天主，因为天主恢复了她的健康。那一夜，那城的居民欢欣鼓舞，以为天主和祂的天使降临到他们中间。\par

16. 他們在那裡住了三天，受到極大的尊敬，生活奢華。之後，他們獲得了旅途所需的糧食，便離開那裡，來到另一個城市。由於該城人口稠密，他們打算在那裡過夜。那城中有一位出眾的婦女：有一次，當她去河邊沐浴時，看哪，可咒罵的撒殫以蛇的形狀跳到她身上，纏繞在她的肚腹上；每當夜晚來臨，他就殘酷地折磨她。這婦女看見了主母夫人瑪利亞，以及她懷中的孩子、主基督，就心中渴望祂，便對主母夫人瑪利亞說：\Quote{主母，請把這孩子給我，好讓我抱抱祂、親親祂。}於是她就把孩子交給那婦女；當孩子被抱到她面前時，撒殫便放開了她，逃離了她，從那日起，那婦女再也沒有看見他。因此，所有在場的人都讚美至高天主，而那婦女也贈予他們豐厚的禮物。\par

17. 次日，那婦女取來香液水為洗滌主耶穌；洗畢，她便用洗過的水，將一部分倒在一個住在那裡、渾身白癩的女孩子身上，為她擦洗。剛一這樣做，女孩的癩病就潔淨了。城裡的人說：\Quote{無疑，若瑟、瑪利亞和那個男孩是神，不是人。}當他們準備離開時，那患癩病的女孩來到他們面前，請求允許她與他們同行。\par

18. 他們應允了她，她就與他們同行。之後，他們來到一座城市，那裡有一位最顯赫的王子的城堡，他設有招待外客的館舍。他們轉入此地；那女孩去見王子的妻子，見她正在哭泣憂傷，便問她為何哭泣。她說：\Quote{不要對我的眼淚感到驚訝；因為我被巨大的痛苦壓倒，至今我還不忍心告訴任何人。}女孩說：\Quote{或許，如果你向我表白並透露，我可能有辦法治療。}公主回答說：\Quote{那麼，請保守這個秘密，不要告訴任何人。我嫁給了這位王子，他是國王，統治許多城市，我與他生活了很久，但他沒有從我得到兒子。當我最終為他生了一個兒子時，他卻患了癩病；他一見到他，就厭惡地轉過身，對我說：\InnerQuote{要麼殺了他，要麼把他交給乳母，帶到一個我們永遠不會再聽到他消息的地方撫養。從此以後，我與你再無關係，我永遠不再見你。}因此，我不知道該怎麼辦，我被悲傷壓倒。唉！我的兒子。唉！我的丈夫。}女孩說：\Quote{我不是說過嗎？我找到了治你病的藥方，我要告訴你。因為我也是個癩病人；但我被天主潔淨了，祂就是耶穌，夫人瑪利亞的兒子。}那婦人問她，她所說的天主在哪裡。女孩說：\Quote{就在你這裡；祂住在同一所房子裡。}她說：\Quote{這怎麼可能？祂在哪裡？}女孩說：\Quote{那裡，若瑟和瑪利亞在那裡；與他們在一起的孩子名叫耶穌；正是祂治好了我的疾病和痛苦。}她說：\Quote{但你是怎樣治好癩病的？你不告訴我嗎？}女孩說：\Quote{為什麼不呢？我從祂母親那裡得到了洗祂的水，澆在自己身上；這樣我就潔淨了癩病。}於是公主起身，邀請他們接受她的款待。她為若瑟在當地人的盛大集會中準備了豐盛的宴會。次日，她取來香液水要洗主耶穌，然後用同樣的水洗了自己的兒子（她把他帶在身邊），她的兒子立刻潔淨了癩病。因此，她歌唱感謝和讚美天主，說：\Quote{有福了，生育你的母親，耶穌啊！你用洗你身體的水潔淨那些與你同性質的人嗎？}此外，她贈予主母夫人瑪利亞厚禮，並以極大的榮耀送她離去。\par

19. 之後，他們來到另一個城市，打算在那裡過夜。於是他們轉入一個新婚男子的家，但他因中巫術而不能與妻子同房；他們與他共度了一夜後，他的束縛就解開了。黎明時分，當他們整裝待發時，新郎不讓他們走，並為他們準備了豐盛的宴會。\par

20. 次日，他们便起程。将近另一座城时，他们看见三个妇女从墓地出来，正在哭泣。\Person{圣母玛利亚}看见她们，便对随行的女子说：问问她们怎么了，或遭遇了什么不幸。那女子问她们，她们却不回答，反问道：你们从哪里来，往哪里去？因为天色已晚，黑夜将至。女子说：我们是旅客，正在寻找可以过夜的客店。她们说：跟我们走吧，今夜与我们同住。于是她们随她们而去，被带入一所新房子，里面装饰华丽，陈设精美。当时正值冬季；那女子进入这些妇女的房间，见她们又在哭泣哀叹。她们旁边站着一头骡子，披着金线织成的鞍褥，面前放着芝麻；妇女们正亲吻它，喂它食物。女子说：我的夫人们，这头骡子是怎么回事？她们含泪答道：你所见的这头骡子，原是我们的兄弟，与我们一母所生。我们的父亲去世后，给我们留下大量财富和这唯一的兄弟。我们尽心为他娶妻，正按人的方式为他准备婚礼。但有些妇人出于相互嫉妒，在我们不知情时对他施了巫术；一夜，天将破晓时，我们屋门紧闭，却见这兄弟变成了骡子，如你现在所见。我们悲痛不已，如你所见，没有父亲安慰我们；世上任何智者、术士或巫师，我们都曾请来，但一无所获。每当心中哀伤难抑，我们就起身带着我们的母亲，到父亲坟前哭泣，然后再回来。\par

21. 那女子听了这些话，便说：放心吧，不要哭；你们灾祸的医治近了，就在你们身边，就在你们家中。因为我本是一个癞病人，但当我看见那妇人和与她同行的幼童，名叫\Person{耶稣}，我便用祂母亲给祂洗澡的水洒在我身上，就痊愈了。我知道\Person{主耶稣}也能医治你们的痛苦。起来，到我的主母\Person{玛利亚}那里去；带她到你们屋里，向她倾诉你们的秘密；恳求她怜悯你们。妇女们听了女子的话，急忙来到\Person{圣母玛利亚}面前，带她进了她们的房间，坐在她面前哭泣，说：我们的主母，\Person{圣母玛利亚}，可怜你的婢女吧！因为我们没有长辈，也没有家主——既无父亲，也无兄弟——与我们同住；但你看见的这头骡子原是我们的兄弟，是妇女们用巫术把它变成这样的。因此我们恳求你怜悯我们。\Person{圣母玛利亚}因她们的遭遇而忧伤，便抱起\Person{主耶稣}，放在骡背上；她也与妇女们一同流泪，对\Person{耶稣基督}说：唉，我的儿子，求你用你的大能治愈这头骡子，使它恢复理智，如从前一样。\Person{圣母玛利亚}说完这话，骡子的形貌就变了，成了一名青年，毫无瑕疵。于是他和他的母亲、姊妹们朝拜\Person{圣母玛利亚}，把那孩子高举过顶，开始亲吻祂，说：\Quote{那生育你的是有福的，哦\Person{耶稣}，世界的救主；那得见你福分的眼睛是有福的。}\par

22. 姊妹二人又对她们的母亲说：我们的兄弟，因\Person{主耶稣基督}的帮助，并因这女子的有益指引——她向我们指出了\Person{玛利亚}和她的儿子——已恢复人形。如今，既然我们的兄弟未婚，将他们的婢女这女子许配给他为妻，倒很合适。于是她们请求\Person{圣母玛利亚}，得到同意后，为那女子举行了盛大的婚礼；她们的忧愁变为喜乐，捶胸变为跳舞；她们开始欢喜、快乐、雀跃、歌唱——因着莫大的喜乐，她们身着最华美的盛装。然后她们开始吟诵赞美诗，说：\Quote{哦\Person{耶稣}，达味之子，你使忧愁变为喜乐，使哀叹变为欢舞！}\Person{若瑟}和\Person{玛利亚}在那里住了十天。此后他们动身，受到这些人的盛情款待；他们向他们告别，含泪返回，尤其那女子更是不舍。\par

23. 他们离开那地，来到一片荒野；听说那里有强盗出没，\Person{若瑟}和\Person{圣母玛利亚}决定连夜穿越此地。他们正走着，忽见两个强盗躺在路上，还有许多同伙睡在他们旁边。那两个抓住他们的强盗，名叫\Person{提图斯}和\Person{杜马库斯}。\Person{提图斯}于是对\Person{杜马库斯}说：我恳求你让这些人自由过去，免得我们的同伴看见他们。\Person{杜马库斯}不肯，\Person{提图斯}又对他说：你从我这里拿去四十个银币，以此为质。同时他把腰间的带子递给他作抵押，免得他开口说话。\Person{圣母玛利亚}见那强盗善待他们，就对他说：\Quote{上主天主必以右手扶持你，赐你罪过得赦。}\Person{主耶稣}回答，对祂的母亲说：\Quote{三十年后，我的母亲，犹太人要在耶路撒冷把我钉在十字架上，这两个强盗也要与我同钉，\Person{提图斯}在我右边，\Person{杜马库斯}在我左边；那日之后，\Person{提图斯}要在我之前进入天堂。}她说：\Quote{我的儿子，但愿天主保佑你免于此难。}他们从那里往一座偶像之城去；将近时，那城变为沙丘。\par

24. 他们于是转向那棵西克莫无花果树，即今所称\Place{马塔雷亚}；\Person{主耶稣}在\Place{马塔雷亚}使一泉水涌出，\Person{圣母玛利亚}在其中洗了祂的衬衫。因\Person{主耶稣}的汗水洒在那里，那地方便出产香液。\par

25. 他们从那里下到门菲斯，见到法老，在埃及住了三年；主耶稣在埃及行了极多奇迹，这些奇迹既未记于\WorkTitle{童年福音}，也未记于\WorkTitle{正典福音}。\par

26. 三年结束时，他离开埃及，返回了。当他们到达犹太时，若瑟害怕进入；但听说黑落德已死，他的儿子阿赫劳斯继位，他虽害怕，还是进入了犹太。主的一位天使显现给他，说：\Quote{若瑟，去纳匝肋城，住在那里。}\par

实在是奇妙，世界的主竟这样被抱着在世界各地行走！\par

27. 此后，他们进入白冷城，看见那里许多孩子患了严重的眼疾，并因此濒临死亡。有一个妇人带着一个病重的儿子，他已近死亡，她把他带到圣母玛利亚跟前，那时圣母正在洗耶稣基督。妇人对她说：\Quote{圣母玛利亚，请看看我的这个儿子，他正遭受重病之苦。}圣母玛利亚听了，说：\Quote{取一点我洗我儿子的水，洒在他身上。}妇人便照圣母玛利亚的话取了点水，洒在她儿子身上。做毕，他的病情减轻了；睡了一会儿后，他平安无恙地醒来。他的母亲为此喜乐，再次把他带到圣母玛利亚面前。圣母对她说：\Quote{感谢天主，因为他已治愈了你的这个儿子。}\par

28. 在同一地方有另一个妇人，是那位儿子刚痊愈的妇人的邻居。她的儿子也罹患同样的疾病，眼睛几乎失明，她日夜哭泣。已痊愈孩子的母亲对她说：\Quote{你为什么不把你的儿子带到圣母玛利亚那里，就像我当初带我的儿子（他几乎死了）那样？他用洗过她儿子耶稣身体的水就好了。}那妇人听了，也去取了同样的水，给她儿子洗了，他的身体和眼睛立刻痊愈了。当她带儿子到圣母玛利亚那里，并告知所发生的一切时，圣母玛利亚吩咐她为儿子恢复健康而感谢天主，并且不要告诉任何人这件事。\par

29. 在同一城中有两个妇人，同为一个男子的妻室，各有一个儿子患热病。一个名叫玛利亚，她的儿子名叫克娄帕。她起身抱起儿子，走到耶稣的母亲圣母玛利亚那里，献上一件美丽的斗篷，说：\Quote{圣母玛利亚，请收下这件斗篷，并给我一条小绷带。}圣母玛利亚照做了；克娄帕的母亲离去，用它做了一件衬衫，给儿子穿上。于是他的病痊愈了；但她对手的儿子却死了。从此她们之间产生了仇恨。她们轮流做家务；轮到克娄帕的母亲玛利亚时，她热炉子烤面包；去取揉好的面团时，她把儿子克娄帕留在炉旁。她的对手见他独自一人——炉子因火燃烧而极热——抓住他扔进炉里，然后自己走开了。玛利亚回来，见儿子克娄帕躺在炉中欢笑，炉子完全冷却，好像从未有火靠近过，就知道是她的对手把他扔进了火里。于是将他拉出，带到圣母玛利亚那里，告诉她所发生的事。圣母说：\Quote{保持沉默，不要告诉任何人这件事；因为我担心你若泄露会遭遇不测。}此后，她的对手去井边打水；见克娄帕在井边玩耍，附近无人，便抓住他扔进井里，然后回家。一些去井边打水的人看见那男孩坐在水面上；于是下去把他拉上来。他们对那男孩大为惊叹，赞美天主。随后他的母亲来了，抱起他，哭着到圣母玛利亚那里，说：\Quote{我的圣母，请看我的对手对我的儿子做了什么，她如何把他扔进井里；她迟早会杀害他的。}圣母玛利亚对她说：\Quote{天主会为你报复她。}此后，当她的对手去井边打水时，脚被绳子缠住，掉进了井里。有人来拉她上来，却发现她头骨碎裂，骨头折断。这样她凄惨地死去；在她身上应验了那句话：\Quote{他们挖了深井，却掉进了自己所预备的坑中。}\par

30. 另有一个妇人有双生子染病，其中一个死了，另一个奄奄一息。母亲哭泣着将他抱起，带到圣母玛利亚那里，说：\Quote{我的圣母，请帮助我，扶持我。因为我本有两个儿子，一个刚刚埋葬，另一个气息奄奄。看，我要如何恳求并祈祷天主。}她开始说：\Quote{主啊，你是慈悲、怜悯、充满慈爱的。你赐给我两个儿子，你已取走一个；至少把这一个留给我。}于是圣母玛利亚见她痛哭的热忱，怜悯了她，说：\Quote{把你的儿子放在我儿子的床上，用他的衣服盖着他。}当她把儿子放在基督躺卧的床上时，他已经闭目死去；但主耶稣基督衣服的气味一传到那男孩身上，他睁开眼睛，大声呼唤母亲，要面包，取过来吸吮。然后他的母亲说：\Quote{圣母玛利亚，现在我知道天主的德能住在你内，以致你的儿子医治那些与他有相同本性的人，只要他们一摸他的衣服。}这个被治愈的男孩就是在福音中被称为巴尔多禄茂的。\par

31. 此外，那里有一个患麻风病的妇人，她来到\Person{玛利亚}夫人——\Person{耶稣}的母亲——跟前，说：夫人，请帮助我。\Person{玛利亚}夫人回答说：你寻求什么帮助？是金子还是银子？还是要你的身体从麻风病中得到洁净？那妇人问：谁能赐我这事？\Person{玛利亚}夫人对她说：稍等片刻，等我洗了我的儿子\Person{耶稣}，并把他放到床上。那妇人照\Person{玛利亚}吩咐的等候；当她把\Person{耶稣}放到床上后，便递给妇人洗过祂身体的水，说：取一点这水，倒在你的身上。她刚一这样做，就洁净了，于是赞美感谢天主。\par

32. 因此，那妇人与她同住了三天后离去；来到一座城，看见那里有一位首领，他已娶了另一位首领的女儿。但当那首领看见这妇人时，发现她两眼之间有麻风病的标记，形状像一颗星；于是婚姻被解除，成了无效。那妇人看见他们如此悲伤、痛哭，便问他们忧伤的原因。但他们说：不要问我们的情形，因为我们的忧伤无法对任何活人诉说，只能对我们自己透露。然而她坚持请求，恳求他们告诉她，说她或许能告诉他们一种疗法。当他们把那女子指给她看，以及她两眼之间出现的麻风病迹象时，那妇人一看见就说：你们在这里看见的我，也曾患同样的病，那时因有件事碰巧到了\Place{白冷}。在那里进了一个山洞，看见一位名叫\Person{玛利亚}的妇人，她的儿子名叫\Person{耶稣}；当她看见我是麻风病人时，就可怜了我，递给我洗过她儿子身体的水。我用那水洒了身体，就出来洁净了。于是那妇人对她说：夫人，你岂不愿意起来与我们同去，领我们见\Person{玛利亚}夫人吗？她答应了；她们起来，带着贵重的礼物到\Place{白冷}去了。她们进去，献上礼物，把带来的患麻风病的女子指给她看。\Person{玛利亚}夫人就说：愿主\Person{耶稣基督}的怜悯降临于你；又把一点洗过\Person{耶稣基督}身体的水交给她们，吩咐给那可怜的妇人洗浴。这事完成后，她立刻痊愈了；她们和所有在场的人都赞美天主。于是她们欢欢喜喜地回到自己的城里，赞美主所行的事。当那首领听说他的妻子痊愈了，就把她接回家，举行了第二次婚礼，并为妻子恢复健康感谢天主。\par

33. 那里还有一个被\Person{撒殚}折磨的年轻妇人；因为那可恶的魔鬼常以一条巨龙的形象出现在她面前，准备吞吃她。他还吸干了她所有的血，使她像一具死尸。每当它走近她时，她就双手抱头，哭喊道：苦啊，苦啊我，没有人能救我脱离那可恶的巨龙！她的父亲、母亲，以及所有在她身边或看见她的人，都为她的命运哀哭；人们围着她，聚成一群，都哭泣哀叹，尤其是当她哭着说：啊，我的兄弟和朋友，难道没有人能救我脱离那个凶手吗？那被治愈了麻风病的首领的女儿，听到那女子的声音，就上到她城堡的屋顶，看见她双手抱头哭泣，所有人群也都围着她哭泣。于是她问那附魔妇人的丈夫，他妻子的母亲是否还在世。他回答说她的父母都健在，她就说：派人叫她母亲到我这里来。她看见他派人去叫，那母亲来了，就说：那个疯癫的女儿是你的女儿吗？是的，夫人，那悲伤哭泣的妇人说，她是我的女儿。首领的女儿回答说：请保守我的秘密，因为我向你承认，我从前是个麻风病人；但现在\Person{玛利亚}夫人，\Person{耶稣基督}的母亲，已经医治了我。但如果你希望你的女儿痊愈，就带她到\Place{白冷}去，寻找\Person{耶稣}的母亲\Person{玛利亚}，相信你的女儿会痊愈；我确实相信你会带着痊愈的女儿欢欢喜喜地回来。那妇人一听首领女儿的话，就急忙领走她的女儿；到了所指的地方，去见\Person{玛利亚}夫人，向她透露她女儿的情况。\Person{玛利亚}夫人听了她的话，就给了她一点洗过她儿子\Person{耶稣}身体的水，吩咐她倒在她女儿身上。她还从主\Person{耶稣}的衣服上给了一条襁褓布，说：拿这块布，每次看见你的仇敌时就把它展现在它面前。然后她问候了她们，打发她们走了。\par

34. 于是，她们离开她，回到自己的地区，当\Person{撒殚}惯常攻击她的时辰临近时，那可恶的魔鬼就以巨龙的形象出现在她面前，女子看见它就害怕。她母亲对她说：我女儿，不要怕；让它走近你，然后把\Person{玛利亚}夫人给我们的那块布给它看，让我们看看会发生什么。\Person{撒殚}于是以可怕的巨龙形象走近，女子的身体因害怕它而颤抖；但她刚拿出那块布，放在头上，用布遮住眼睛，就有火焰和火炭从中射出，投向巨龙。哦，何等伟大的奇迹！巨龙一看见主\Person{耶稣}的布，火焰从中射出，投向它的头和眼睛，它就大声喊叫：\Person{耶稣}，\Person{玛利亚}之子，我与你有何相干？我该逃到哪里躲避你？它极其恐惧地转身离开那女子，此后再也没有在她面前出现。那女子从此摆脱了它，赞美感谢天主，所有在场目睹那奇迹的人也一同赞美感谢。\par

35. 另有一妇人，居于同地，其子为撒殚所折磨。此子名犹达斯，每当撒殚附身，辄咬近旁之人；若近处无人，则咬己手及他肢体。此可怜之母闻圣母玛利亚及其子耶稣之声名，遂起，携其子犹达斯至圣母前。其时，雅各伯与若瑟已抱主耶稣童婴，与诸童嬉戏；他们出户坐定，主耶稣亦在焉。附魔之犹达斯前来，坐于耶稣右侧；忽为撒殚如前所侵，欲咬主耶稣，却不能；然击耶稣右肋，耶稣遂啼泣。撒殚即时自该童而出，如狂犬奔逃。此击耶稣之童，撒殚以犬形出之者，乃依斯加略人犹达斯，后出卖耶稣于犹太人；犹太人亦以长枪刺透耶稣之右肋，即犹达斯所击之处。\BibleRef{若 19:34}\par

36. 主耶稣自诞生满七岁时，一日与同龄之童嬉戏。彼等以泥土塑驴、牛、鸟及其它动物之像，各夸己技，自赞其工。主耶稣谓诸童曰：\Quote{我所塑之像，我命之行走。}诸童问曰：\Quote{尔岂造物主之子乎？}主耶稣命之行走，像即跳跃；又命之止，像即静止。又塑鸟雀之像，命飞即飞，命止即止，授之食饮，即食即饮。诸童归告父母，其父曰：\Quote{我儿，慎勿复与彼游，彼乃巫师；远之避之，此后勿与戏。}\par

37. 一日，主耶稣奔跑嬉戏，过染匠撒冷之店；店中有布匹甚多，待染。主耶稣入店，尽取诸布，投于盛满靛蓝之缸中。撒冷归见布毁，大声呼叱耶稣曰：\Quote{玛利亚之子，尔何为此？尔辱我于众同乡前！盖各人欲求合己之色，尔竟来尽毁之。}主耶稣答曰：\Quote{凡尔欲改之色，我必为尔易之。}随即自缸中取出布匹，各如染匠所欲之色，尽取而出。犹太人见此奇迹奇事，遂赞美天主。\par

38. 若瑟常巡行全城，携主耶稣同往；人召若瑟作门、奶桶、床、柜等木工，主耶稣随往。凡若须造之物，或长或短，或广或狭，主耶稣伸手抚之，即成若瑟所愿。若瑟无需亲手劳作，盖其木工技艺未精。\par

39. 一日，耶路撒冷王召若瑟曰：\Quote{我欲尔为我造一宝座，合于我所常坐之处。}若瑟遵命，即日兴工，在王宫二年，始成宝座。及抬至其所，见每边较规定之尺寸短两掌。王见之怒，若瑟惧王甚，终夜不食，未尝一物。主耶稣问其何惧，若瑟曰：\Quote{我二年之功尽废矣。}主耶稣曰：\Quote{勿惧，勿灰心；尔持宝座一边，我持另一边，可修整之。}若瑟如言，各持一边，宝座遂正，适符所量。旁观者见此奇迹，惊异赞美天主。宝座所用之木，乃撒罗满达味之子时所称道者，即多种多样之木。\par

40. 另一日，主耶稣出行于路，见诸童集嬉，遂从之；诸童却避匿。主耶稣至一房门，见数妇立彼，问童何往。妇答曰：\Quote{无人。}耶稣又曰：\Quote{尔等所见炉中者谁？}妇曰：\Quote{三岁之山羊羔。}主耶稣呼曰：\Quote{山羊羔，出来就尔牧人。}诸童即化山羊羔之形而出，环舞于耶稣四周。妇人大惊，战栗，速求耶稣，朝拜曰：\Quote{我主耶稣，玛利亚之子，尔诚以色列之善牧；求尔垂怜尔婢，今立尔前，未尝怀疑；盖我主来为医治，非为毁灭。}主耶稣答曰：\Quote{以色列子民在万民中，如厄提约丕雅人然。}妇人曰：\Quote{主，尔知万事，无隐于尔；今恳求尔，以尔慈爱，复此诸童本形。}主耶稣遂曰：\Quote{众童，来，我侪游戏。}言毕，诸羊羔立化童形，诸妇在旁目睹。\par

41. 时在阿达尔月，耶稣效仿君王之仪，聚集众童。他们将衣服铺在地上，耶稣便坐于其上。又以花冠加于其首，如内侍般侍立于其左右，仿佛祂是君王。凡路过者，童众强拉之，说：\Quote{前来朝拜君王，然后离去。}\par

42. 其间，有人抬着一男孩前来。原来此男孩与同龄人上山拾柴，见一鹧鸪巢；伸手取蛋时，巢中毒蛇啮之，以致他呼救。同伴急往，见其卧地如死。亲属遂来抬他回城。行至主耶稣如王坐处，众童如仆环立，童众急迎被蛇咬者，对其亲属说：\Quote{来朝拜君王。}他们因忧愁不愿前往，童众强行拉之。及至主耶稣前，主问为何抬此男孩。答曰毒蛇所咬。主耶稣对童众说：\Quote{我们去杀那蛇。}男孩父母请求离去，因其子濒死；童众答曰：\Quote{尔等未闻君王言\InnerQuote{我们去杀蛇}乎？岂可不从？}遂违其意，抬床返回。至巢穴，主耶稣问童众：\Quote{此是蛇穴吗？}答曰是；蛇闻主召，立时出来，向祂顺服。主说：\Quote{去吸尽你注入此孩之毒。}蛇遂爬向男孩，吸尽其毒。然后主耶稣咒诅蛇，蛇立时裂开；主耶稣用手抚慰男孩，他便痊愈。男孩哭泣，耶稣说：\Quote{勿哭，不久你将作我门徒。}此即西满加纳人，福音中所提及者。\par

43. 另一日，若瑟遣其子雅各伯拾柴，主耶稣与之同往。至柴薪处，雅各伯始拾柴，忽有毒蛇啮其手，遂哭号。主耶稣见之，上前向伤口吹气，即刻痊愈。\par

44. 一日，主耶稣又与童众在屋顶玩耍，一童从高处坠下，当即气绝。众童四散，惟主耶稣独留屋顶。男孩亲属前来对主耶稣说：\Quote{是你将我儿从屋顶推下。}主否认，他们喊道：\Quote{我儿死了，凶手在此。}主耶稣对他们说：\Quote{莫妄加罪名；若不信，可问男孩本人，以明真相。}主耶稣遂下，立于尸旁，大声说：\Quote{则诺，则诺，谁将你从屋顶推下？}死童答道：\Quote{我主，非你推我，乃某人将我掷下。}主命旁观者听其言，众人皆因这奇迹赞美天主。\par

45. 昔时，圣母玛利亚命主耶稣去井边取水。及至取水，已满之瓶撞物而碎。主耶稣展开手帕，收水携回给母亲；母亲惊异，将所见一切默存心中。\par

46. 又一日，主耶稣与童众在溪边玩耍，又筑小鱼池。主耶稣造了十二只麻雀，置于池周，每边三只。时值安息日。有一犹太人，哈难之子，前来见其作为，怒斥道：\Quote{安息日你们造泥像吗？}遂急毁其鱼池。主耶稣却对其所造之雀拍手，雀即飞去，啾啾而鸣。\par

哈难之子又至耶稣的鱼池，以鞋踢之，池水尽失。主耶稣对他说：\Quote{如水消失，你的性命亦将如是消失。}那孩子即刻枯干。\par

47. 另有一次，主耶稣晚间与若瑟回家，遇一男孩猛撞祂，致其跌倒。主耶稣对他说：\Quote{你推倒我，也必跌倒不起。}顷刻间，男孩倒地而亡。\par

48. 此外，在耶路撒冷有一个人，名叫\Person{匝凯}，他教男孩们。他对\Person{若瑟}说：\Quote{啊，\Person{若瑟}，你为什么不带耶稣来学习字母？}\Person{若瑟}同意了，并告知了\Person{圣母玛利亚}。于是他们带他去见老师；老师一看见他，就给他写出字母表，叫他念\OriginalTerm{阿勒弗}{Aleph}。当他念了\OriginalTerm{阿勒弗}{Aleph}后，老师吩咐他念\OriginalTerm{贝特}{Beth}。主耶稣对他说：\Quote{请先告诉我字母\OriginalTerm{阿勒弗}{Aleph}的意思，然后我再念\OriginalTerm{贝特}{Beth}。}当老师威胁要鞭打他时，主耶稣向他解释了\OriginalTerm{阿勒弗}{Aleph}和\OriginalTerm{贝特}{Beth}字母的意思；也解释了哪些字母形状是直的，哪些是弯的，哪些卷成螺旋形，哪些有点，哪些没有，为什么一个字母在另一个之前；还有许多其他事情他开始述说并阐明，这些老师自己从未在任何书中听过或读过。主耶稣又对老师说：\Quote{听着，我要念给你听。}于是他清晰明确地重复\OriginalTerm{阿勒弗}{Aleph}、\OriginalTerm{贝特}{Beth}、\OriginalTerm{基默尔}{Gimel}、\OriginalTerm{达肋特}{Daleth}，直到\OriginalTerm{陶}{Tau}。老师惊讶地说：\Quote{我想这孩子生在\Person{诺厄}之前。}他转身对\Person{若瑟}说：\Quote{你给我带来一个比所有老师都更有学问的孩子来受教。}又对\Person{圣母玛利亚}说：\Quote{你的这个儿子不需要教导。}\par

49. 此后，他们带他去见另一位更有学问的老师；那老师一看见他，就说：\Quote{念\OriginalTerm{阿勒弗}{Aleph}。}当他念了\OriginalTerm{阿勒弗}{Aleph}后，老师吩咐他念\OriginalTerm{贝特}{Beth}。主耶稣回答他说：\Quote{请先告诉我字母\OriginalTerm{阿勒弗}{Aleph}的意思，然后我再念\OriginalTerm{贝特}{Beth}。}当老师这时举起手鞭打他时，他的手立刻枯干，他就死了。然后\Person{若瑟}对\Person{圣母玛利亚}说：\Quote{从今以后，我们不要让他出门，因为凡反对他的人都必遭击杀而死。}\par

50. 当他十二岁时，他们带他到\Place{耶路撒冷}过庆节。庆节结束后，他们回去了；但主耶稣留在圣殿中，在以色列子民的教师、长老和经师中间，向他们提出各种科学问题，并轮到他作答。因为他对他们说：\Quote{\OriginalTerm{默西亚}{Messias}是谁的儿子？}他们回答他说：\Quote{达味的儿子。}既然如此，他说：\Quote{达味在圣神里怎么称他为主，当他说：\InnerQuote{上主对我主说：你坐在我右边，等我使你的仇敌作你的脚凳？}}教师的领袖又问他说：\Quote{你读过这些书吗？}主耶稣说：\Quote{书和书中所包含的都读过。}他解释了书、法律、诫命、典章、以及先知书中所含的奥秘——这些是任何受造物的理解力所不能及的。因此，那教师说：\Quote{我至今没有达到也没有听说过这样的知识：请问，你认为这孩子将会是谁？}\par

51. 当时在场的一位哲学家，一位精通天文学者，问主耶稣是否研究过天文学。主耶稣回答了他，并解释了天球和天体的数目、它们的本性和运行；它们冲的位置；它们的相位：\OriginalTerm{三分相}{triangular}、\OriginalTerm{四分相}{square}和\OriginalTerm{六分相}{sextile}；它们的运行：顺行和逆行；二十四分之一和二十四分之一的六十分之一；以及其它超出理智范围的事。\par

52. 在那些哲学家中，还有一位非常精通自然科学的，他问主耶稣是否研究过医学。主耶稣回答了他，解释了\OriginalTerm{物理学}{physics}和\OriginalTerm{形而上学}{metaphysics}、\OriginalTerm{超物理学}{hyperphysics}和\OriginalTerm{下物理学}{hypophysics}，同样也解释了身体的力量和\OriginalTerm{体液}{humours}，以及它们的作用；也解释了肢体和骨骼、静脉、动脉和神经的数目；也解释了热和干、冷和湿的作用，以及它们所产生的结果；灵魂对身体的作用，及其知觉和能力；说话、愤怒、欲望的能力作用；最后，它们的结合与分离，以及其他任何受造理智所不能及的事。于是那哲学家站起来，朝拜了主耶稣，说：\Quote{主啊，从今以后，我要做你的门徒和奴仆。}\par

53. 当他们彼此谈论这些和其他事情时，\Person{圣母玛利亚}来了，她和\Person{若瑟}四处寻找他已有三天。因此她看见他坐在教师中间，向他们提问，并轮到他回答，就对他说：\Quote{我的儿子，你为什么这样对待我们？看，你的父亲和我很辛苦地找你。}但他说：\Quote{你们为什么找我？你们不知道我必须在我父的家里吗？}但他们不明白他对他们说的话。然后那些教师问\Person{玛利亚}他是否是她的儿子；当她表示确实如此时，他们说：\Quote{玛利亚，你是有福的，因为你生了这样一个儿子。}于是他们一起回到\Place{纳匝肋}，他事事服从他们。他的母亲把这些话都默存在心里。主耶稣在智慧和身量上，并在天主和人前的恩宠上，渐渐地增长。\BibleRef{路 2:46}\par

54. 从那天起，他开始隐藏他的奇迹、奥秘和秘密，并专心致志于法律，直到他满了三十岁，那时圣父在\Place{约旦河}公开宣告他，从天上降下这样的声音：\Quote{这是我的爱子，我所喜悦的；}圣神以白鸽的形象临在。\par

55. 这就是我们以恳求朝拜的那一位，他赐予我们存在和生命，并使我们从母腹中出生；他为了我们取了人的身体，救赎了我们，为要以永恒的慈悲拥抱我们，并按照他的慷慨、恩惠、宽仁和慈爱向我们显示他的怜悯。愿荣耀、恩惠、权能和统治归于他，从今时直到永远。阿们。\par

\WorkTitle{耶稣童年福音}到此结束。仰赖至高天主的助佑，依照我们在原稿中所发现的。\par

\SourceRecord{Translated by Alexander Walker.；Ante-Nicene Fathers；Vol. 8.；Edited by Alexander Roberts, James Donaldson, and A. Cleveland Coxe.；Buffalo, NY: Christian Literature Publishing Co.,；1886.；<http://www.newadvent.org/fathers/0806.htm>.}
